% =====================================================================
\chapter{Cuantificadores y negación}
\label{tema:2}
% =====================================================================

Casi ningún teorema habla de un único objeto: afirman algo de \emph{todos} los enteros, de \emph{alguna} raíz de un polinomio o de que \emph{para cada} $\varepsilon$ \emph{existe} un $\delta$. Este tema introduce los predicados y los cuantificadores, estudia el efecto del orden de los cuantificadores y da reglas sistemáticas para negar enunciados, que serán imprescindibles en las demostraciones por contrarrecíproco y por reducción al absurdo. Se siguen las §2.7–2.10 de~\cite{hammack}, el capítulo~2 de~\cite{velleman} (§2.1 y §2.2), la §3.6 de~\cite{mcs} y la guía de negación del curso CS103~\cite{cs103}.

\section{Predicados}

\begin{definicion}[Predicado]\label{def:predicado}
Sea $U$ un conjunto no vacío, llamado \emph{universo} o dominio. Un \emph{predicado} sobre $U$ es un enunciado $P(x)$ que contiene una variable $x$ y que se convierte en una proposición cada vez que se sustituye $x$ por un elemento concreto de $U$. El \emph{conjunto de verdad} de $P$ es $\{x\in U: P(x)\text{ es verdadera}\}$. Del mismo modo se definen predicados de varias variables, $P(x,y)$, con $(x,y)\in U\times W$.
\end{definicion}

\addtocounter{definicion}{1}%

\section{Cuantificador universal y existencial}

\begin{definicion}[Cuantificadores]\label{def:cuantificadores}
Sea $P(x)$ un predicado sobre $U$.
\begin{enumerate}
  \item $\forall x\in U\ P(x)$ («para todo $x$ de $U$, $P(x)$») es verdadera si $P(a)$ es verdadera para cada $a\in U$.
  \item $\exists x\in U\ P(x)$ («existe $x$ en $U$ tal que $P(x)$») es verdadera si $P(a)$ es verdadera para al menos un $a\in U$.
\end{enumerate}
Cuando el universo está claro se escribe $\forall x\,P(x)$ y $\exists x\,P(x)$. La variable $x$ está \emph{ligada} por el cuantificador: $\forall x\,P(x)$ y $\forall t\,P(t)$ dicen lo mismo.
\end{definicion}

Los cuantificadores se combinan con las conectivas. Las formas más frecuentes son la \emph{universal condicional}, $\forall x\,\bigl(P(x)\Rightarrow Q(x)\bigr)$, que traduce «todo $P$ es $Q$», y la \emph{existencial conjuntiva}, $\exists x\,\bigl(P(x)\wedge Q(x)\bigr)$, que traduce «algún $P$ es $Q$». Se abrevian como $\forall x\in A\ Q(x)$ y $\exists x\in A\ Q(x)$ cuando $P(x)$ es «$x\in A$».

\addtocounter{definicion}{1}%

\begin{observacion}
Si $A=\varnothing$, el enunciado $\forall x\in A\ Q(x)$ es verdadero, sea cual sea $Q$: no hay ningún $x\in A$ que lo incumpla. Es la versión cuantificada del condicional vacuamente verdadero del \cref{tema:1}. En cambio, $\exists x\in\varnothing\ Q(x)$ es siempre falso.
\end{observacion}

\section{Cuantificadores anidados}

Con predicados de varias variables, los cuantificadores se anidan. El orden importa cuando se mezclan cuantificadores de distinto tipo.

\addtocounter{definicion}{1}%

\begin{proposicion}\label{prop:orden}
Para todo predicado $P(x,y)$,
\[
\exists y\,\forall x\,P(x,y)\ \Longrightarrow\ \forall x\,\exists y\,P(x,y),
\]
y la implicación recíproca es falsa en general.
\end{proposicion}

Cuantificadores del mismo tipo sí pueden intercambiarse: $\forall x\,\forall y\,P(x,y)\eqlog\forall y\,\forall x\,P(x,y)$ y $\exists x\,\exists y\,P(x,y)\eqlog\exists y\,\exists x\,P(x,y)$.

\addtocounter{definicion}{1}%

\section{Negación de enunciados cuantificados}

\begin{teorema}[Negación de cuantificadores]\label{teo:negacion}
Para todo predicado $P(x)$ sobre un universo $U$,
\[
\neg\,\forall x\,P(x)\ \eqlog\ \exists x\,\neg P(x),\qquad
\neg\,\exists x\,P(x)\ \eqlog\ \forall x\,\neg P(x).
\]
\end{teorema}

El \cref{teo:negacion} y las leyes de la \cref{prop:leyes} dan un procedimiento mecánico: la negación «atraviesa» los cuantificadores cambiando $\forall$ por $\exists$ y viceversa, y al llegar a la parte proposicional se aplican De Morgan y la negación del condicional, $\neg(P\Rightarrow Q)\eqlog P\wedge\neg Q$. Con universos acotados, el cuantificador conserva su dominio: $\neg\,\forall x\in A\ Q(x)\eqlog\exists x\in A\ \neg Q(x)$.

\addtocounter{definicion}{1}%

\section{Traducción y unicidad}

Las expresiones del lenguaje natural ocultan a menudo los cuantificadores. «Los primos mayores que~2 son impares» es universal; «un cuadrado puede ser par» es existencial; «solo los múltiplos de~6 son múltiplos de~2 y de~3» significa «si $n$ es múltiplo de~2 y de~3, entonces es múltiplo de~6». Antes de demostrar un enunciado conviene escribirlo con cuantificadores explícitos~\cite[§2.9]{hammack}.

\begin{definicion}[Existencia y unicidad]\label{def:unicidad}
$\exists!\,x\,P(x)$ («existe un único $x$ tal que $P(x)$») abrevia
\[
\exists x\,\Bigl(P(x)\wedge\forall y\,\bigl(P(y)\Rightarrow y=x\bigr)\Bigr).
\]
\end{definicion}

Así, $\exists!$ contiene dos afirmaciones: hay al menos un elemento con la propiedad (\emph{existencia}) y dos elementos con la propiedad son iguales (\emph{unicidad}). El \cref{tema:6} explica cómo se demuestra cada parte.

\addtocounter{definicion}{1}%

\addtocounter{definicion}{1}%

